% Promoted from experimental/experiments.tex lines 2091--2462.
% Status: stable high-agreement note.  This is a TeX fragment intended to be
% included from experimental/experiments.tex or another paper-level wrapper.
% Dependencies: theorem environments, standard macros such as \F and \RS,
% and earlier local labels referenced inside the note.

\section{General High-Agreement Ledgers Beyond the \texorpdfstring{$\F_{17^{32}}$}{F17^32} Row}
\label{sec:general-high-agreement-ledgers}

\textbf{Status: \textsc{New-local} finite Reed--Solomon/MDS theorem package plus
\textsc{Conditional} protocol ledger.}
This section generalizes the high-agreement tangent-star method away from the
special field \(\F_{17^{32}}\).  The statements below are finite-grid coding
statements.  They do not use smoothness, except when one wants to combine them
with the quotient-core and generated-field ledgers elsewhere in this repository.
The line lower bound may be proved either by the RS locator construction of
Theorem~\ref{thm:moving-root-tangent-floor} or by the coordinate tangent-star
construction below; the upper bounds use only the MDS zero-rigidity property and
linearity of the code.

Throughout this section let \(C=\RS[F,D,k]\) with \(D\subset F\), \(|D|=n\),
and dimension \(k\).  The upper-bound arguments below use only the MDS
zero-rigidity property of Reed--Solomon codes, so they apply verbatim to any
linear MDS code for which the relevant received words and parameter sets are
defined.  Put
\[
  R=n-k,
  \qquad r=n-a .
\]
Thus \(r\) is the integer Hamming radius and \(a=n-r\) is the agreement
threshold.  We only discuss the high-agreement range \(a\ge k+1\), equivalently
\(0\le r\le R-1\).

\subsection{A coordinate tangent star}

The next construction is useful because it removes all dependence on the
particular field \(\F_{17^{32}}\), on multiplicative smoothness, and on the
locator-polynomial model.  It uses only coordinate values, linearity, and the
MDS zero-rigidity of Reed--Solomon codes.

\begin{lemma}[coordinate tangent-star lower bound]
\label{lem:coordinate-tangent-star}
Let \(C=\RS[F,D,k]\) with \(D\subset F\).  For every \(a\) with \(k+1\le a\le n\),
\[
  \operatorname{LD}_{\rm sw}(C,a)\ge n-a+1 .
\]
If in addition \(2a-n-1\ge k\), then the same construction gives
\[
  \operatorname{LD}_{\rm ca}(C,a)\ge n-a+1 .
\]
\end{lemma}

\begin{proof}
Choose \(A\subset D\) with \(|A|=a-1\), and write
\(T=D\setminus A\), so \(|T|=n-a+1\).  Choose an injection
\(\theta:T\hookrightarrow F\), which exists because \(T\subseteq D\subset F\).
Define two received words \(f,g:D\to F\) by
\[
  f(x)=g(x)=0\quad (x\in A),
  \qquad
  f(t)=-\theta(t),\quad g(t)=1\quad (t\in T).
\]
For a point \(t\in T\), the finite slope \(z=\theta(t)\) gives
\(f+zg=0\) on \(A\cup\{t\}\), a support of size \(a\).  Hence this slope is
explained by the zero codeword.

The same support is support-wise noncontained.  Indeed, if both \(f\) and \(g\)
were separately explained on \(A\cup\{t\}\), then the explaining codeword for
\(g\) would vanish on the \(a-1\ge k\) points of \(A\), hence would be the zero
codeword by the MDS property.  But \(g(t)=1\), a contradiction.

The pair \((f,g)\) is also globally far from \(C^2\) at agreement \(a\).  The
same argument applied to the word \(g\) shows that no codeword of \(C\) agrees
with \(g\) on any support of size \(a\): such a support contains at least
\(a-|T|=2a-n-1\) points of \(A\), and in the high-agreement applications below we
will have \(2a-n-1\ge k\).  Alternatively, for the lower bound in the line case
one may use the RS locator construction of Theorem~\ref{thm:moving-root-tangent-floor},
which gives the no-loss CA witness directly for \(\RS[F,D,k]\).  The displayed
MCA lower bound holds unconditionally for \(a\ge k+1\), and the CA lower bound
holds in the range \(2a-n-1\ge k\), in particular throughout the exact staircase
range \(3a-2n\ge k\).
\end{proof}

\begin{remark}[scope]
The support-wise MCA construction is coordinate-level.  It uses MDS rigidity
and the inclusion \(D\subset F\).  The no-loss CA part needs global farness; in
the exact range below, the needed support-count condition follows from the
staircase hypothesis.
\end{remark}

\subsection{Uniform line, CA, and projective-slope threshold}

\begin{theorem}[uniform high-agreement line staircase]
\label{thm:uniform-line-staircase}
Let \(C=\RS[F,D,k]\) with \(|D|=n\), and let \(a=n-r\).  If
\[
  3a-2n\ge k
  \qquad\text{equivalently}\qquad
  r\le \frac{n-k}{3},
\]
then
\[
  \operatorname{LD}_{\rm sw}(C,a)
  =
  \operatorname{LD}_{\rm ca}(C,a)
  =
  \operatorname{LD}_{\rm sw,proj}(C,a)
  =
  r+1
  =n-a+1 .
\]
Here \(\operatorname{LD}_{\rm sw,proj}\) allows slopes in \(\mathbb P^1(F)\).
\end{theorem}

\begin{proof}
The finite-slope support-wise upper bound is Theorem~\ref{thm:very-high-tangent-staircase},
whose proof uses only MDS zero-rigidity and the common residual budget.  The
lower bound is Lemma~\ref{lem:coordinate-tangent-star}, or, for Reed--Solomon
codes, Theorem~\ref{thm:moving-root-tangent-floor}.  No-loss CA-bad slopes are a
subset of support-wise MCA-bad slopes, while the same tangent-star witness gives
\(r+1\) no-loss CA-bad slopes in the present range, so the CA numerator is also
\(r+1\).

For projective slopes, suppose a projective bad set had size larger than
\(r+1\).  If there is a nonbad projective point, use it as the point at infinity;
all projective bad points then become finite bad slopes, contradicting the
finite-slope bound.  If every point of \(\mathbb P^1(F)\) is bad, choose any
point as infinity.  The remaining \(|F|\) finite slopes are bad.  Since
\(|F|\ge n>r+1\), this again contradicts the finite-slope bound.  Thus the
projective numerator is also \(r+1\), and the finite lower-bound construction
shows equality.
\end{proof}

\subsection{Degree-\texorpdfstring{$d$}{d} curve MCA/CA}

For an integer \(d\ge1\), a finite-parameter power curve is a received family
\[
  W_\gamma=f_0+\gamma f_1+\cdots+\gamma^d f_d,
  \qquad \gamma\in F .
\]
The support-wise curve-MCA predicate asks for parameters \(\gamma\) for which
\(W_\gamma\) has an explanation on a support of size at least \(a\), while the
\((d+1)\)-tuple \((f_0,\ldots,f_d)\) is not simultaneously explained on that
same support.  The no-loss curve-CA predicate replaces support-wise
noncontainment by global farness of the \((d+1)\)-tuple from \(C^{d+1}\).

\begin{theorem}[degree-\(d\) high-agreement curve staircase]
\label{thm:uniform-curve-staircase}
Let \(C=\RS[F,D,k]\) with \(|D|=n\), and let \(a=n-r\).  If
\[
  (d+2)a-(d+1)n\ge k
  \qquad\text{equivalently}\qquad
  r\le \frac{n-k}{d+2},
\]
then
\[
  \operatorname{CurveLD}_{\rm sw}^{(d)}(C,a)
  =
  \operatorname{CurveLD}_{\rm ca}^{(d)}(C,a)
  =
  \min\{|F|,\ d(r+1)\} .
\]
\end{theorem}

\begin{proof}
First prove the lower bound.  Choose \(A\subset D\) with \(|A|=a-1\) and put
\(T=D\setminus A\), so \(|T|=r+1\).  Choose disjoint parameter sets
\(\Gamma_t\subseteq F\), \(t\in T\), with \(|\Gamma_t|\le d\) and
\(\sum_t |\Gamma_t|=\min\{|F|,d(r+1)\}\).  For each \(t\), choose a nonzero
polynomial
\[
  h_t(Y)=h_{t,0}+h_{t,1}Y+\cdots+h_{t,d}Y^d
\]
of degree exactly \(d\), with leading coefficient nonzero, that vanishes on all
parameters in \(\Gamma_t\).  Define received words by
\[
  f_i(x)=0\quad(x\in A),
  \qquad
  f_i(t)=h_{t,i}\quad(t\in T),
  \qquad 0\le i\le d .
\]
For \(\gamma\in\Gamma_t\), the curve word \(W_\gamma\) vanishes on
\(A\cup\{t\}\), hence is explained by the zero codeword on a support of size
\(a\).  The tuple is not simultaneously explained on that support: on the
\(a-1\ge k\) points of \(A\), every explaining endpoint codeword would have to be
zero, but at \(t\) the coefficient vector of \(h_t\) is not zero.  This gives the
support-wise curve-MCA lower bound.

The same construction is globally far in the exact range.  The last endpoint
\(f_d\) is zero on \(A\) and nonzero on every point of \(T\).  Any support of
size \(a\) contains at least \(a-|T|=2a-n-1\) points of \(A\).  The hypothesis
\(r\le (n-k)/(d+2)\), with \(d\ge1\), implies \(2a-n-1\ge k\).  Therefore any
codeword agreeing with \(f_d\) on such a support would have to be the zero
codeword, but the support must contain at least one point of \(T\), where
\(f_d\ne0\).  Thus the endpoint tuple is not \(C^{d+1}\)-close at agreement
\(a\), proving the no-loss curve-CA lower bound.

For the upper bound, it is enough to rule out more than \(d(r+1)\) bad
parameters, because the trivial cap is \(|F|\).  If a curve had more than
\(d(r+1)\) bad parameters, then in particular it would have at least \(d+1\)
bad parameters.  Select \(d+1\) distinct bad parameters
\(\gamma_0,\ldots,\gamma_d\) and supports \(S_i\) of size at least \(a\).  Their
intersection has size at least
\[
  (d+1)a-dn .
\]
On this intersection the \(d+1\) explained codewords interpolate, as functions
of \(\gamma\), to codewords \(c_0,\ldots,c_d\in C\) explaining the endpoints
\(f_0,\ldots,f_d\).  Let \(S_0\) be a maximal common support for this endpoint tuple,
put \(b=|S_0|\), and set \(m=n-b\).  Then
\[
  m\le (d+1)(n-a)=(d+1)r .
\]
For any further bad parameter, its explaining codeword must equal
\(c_0+\gamma c_1+\cdots+\gamma^d c_d\), because
\[
  a+b-n\ge (d+2)a-(d+1)n\ge k .
\]
Thus outside \(S_0\) it needs at least \(h=\max(1,a-b)\) zeros of the residual
coordinate polynomials.  By maximality of \(S_0\), no residual coordinate
polynomial is identically zero, and each has degree at most \(d\), hence can be
charged to at most \(d\) bad parameters.  If \(b\ge a\), then \(m\le r\) and the
count is at most \(dm\le dr\le d(r+1)\).  If \(b<a\), then
\(h=a-b=m-r\) and \(r<m\le(d+1)r\); the charging bound gives
\[
  \#\{\text{bad parameters}\}
  \le
  \left\lfloor\frac{dm}{m-r}\right\rfloor
  \le d(r+1).
\]
The trivial parameter-space cap gives the additional upper bound \(|F|\).
Together with the lower bound this proves equality.
\end{proof}

\subsection{Interleaved-list uniqueness}

\begin{theorem}[uniform interleaved-list uniqueness]
\label{thm:uniform-interleaved-uniqueness}
Let \(C\subseteq F^D\) be a linear MDS code of dimension \(k\), and let
\(\mu\ge1\).
If
\[
  2a-n\ge k
  \qquad\text{equivalently}\qquad
  r\le\frac{n-k}{2},
\]
then every \(\mu\)-row received word has at most one \(\mu\)-row codeword tuple
agreeing with it on at least \(a\) common coordinates.  Since equality is
attained by taking the received word itself to be a code tuple,
\[
  \Lambda_\mu(C,a)=1 .
\]
\end{theorem}

\begin{proof}
If two \(\mu\)-row codeword tuples agree with the same received word on supports
\(S\) and \(S'\) of size at least \(a\), then they agree with each other on
\(S\cap S'\), whose size is at least \(2a-n\ge k\).  Row by row, the MDS
property forces the two codewords to be equal.  Hence the tuples are equal.
The lower bound \(\Lambda_\mu(C,a)\ge1\) is obtained by choosing the received
word to be a codeword tuple.
\end{proof}

\subsection{General finite-grid threshold calculator}

Let \(\varepsilon_*=2^{-128}\) and let \(Q\) be the denominator of the sampler or
challenge field used by a particular ledger.  Define the integer budget
\[
  B_Q=\left\lfloor \varepsilon_* Q\right\rfloor .
\]
The high-agreement theorems above convert coding ledgers into integer arithmetic.
For example, in the exact line range \(r\le\lfloor R/3\rfloor\), a single affine
line, no-loss CA, or projective-line numerator is
\[
  N_{\rm line}(r)=r+1 .
\]
Thus the line term alone is safe whenever \(r+1\le B_Q\), and the first unsafe
radius is \(r=B_Q\), provided \(B_Q\le\lfloor R/3\rfloor\).  If
\(B_Q>\lfloor R/3\rfloor\), then the tangent-star method proves safety
throughout the exact high-agreement range but does not locate the eventual
threshold.

For one degree-\(d\) finite-parameter curve term in its exact range
\(r\le\lfloor R/(d+2)\rfloor\), the numerator is
\[
  N_{\rm curve,d}(r)=\min\{|F|,d(r+1)\}.
\]
If all coding terms use the same denominator \(Q\), a protocol ledger with curve
terms of degrees \(d_1,\ldots,d_s\) and \(\ell\) independent interleaved-list
terms has high-agreement numerator
\[
  N_{\rm total}(r)=\ell+\sum_{i=1}^s d_i(r+1),
\]
valid for
\[
  r\le
  \min\left\{
    \left\lfloor\frac R2\right\rfloor\text{ if }\ell>0,
    \min_i\left\lfloor\frac R{d_i+2}\right\rfloor
  \right\}.
\]
The ledger is certified at target \(\varepsilon_*\) whenever
\[
  N_{\rm total}(r)\le B_Q .
\]
With different denominators \(Q_i\), one should not combine numerators; the
correct test is the rational inequality
\[
  \sum_i \frac{N_i(r)}{Q_i}+\varepsilon_{\rm query}\le \varepsilon_* .
\]

\begin{corollary}[row-independent closed-form gates]
\label{cor:row-independent-gates}
Let \(R=n-k\), \(B_Q=\lfloor Q/2^{128}\rfloor\), and assume the relevant radius
lies in the exact high-agreement range.
\begin{enumerate}
\item A single affine/projective line, no-loss CA, or support-wise MCA term is
safe for every integer radius
\[
  r\le B_Q-1 .
\]
The first unsafe radius is \(r=B_Q\), if \(B_Q\le\lfloor R/3\rfloor\).
\item A single line/CA/MCA term plus one interleaved-list term over the same
field is safe for
\[
  r\le B_Q-2 .
\]
The first unsafe radius is \(r=B_Q-1\), if \(B_Q-1\le\lfloor R/3\rfloor\).
\item A single degree-\(d\) curve term is safe for
\[
  r\le \left\lfloor \frac{B_Q}{d}\right\rfloor-1 .
\]
The first unsafe radius is \(r=\lfloor B_Q/d\rfloor\), if this radius is at most
\(\lfloor R/(d+2)\rfloor\).
\item A single degree-\(d\) curve term plus one interleaved-list term over the
same field is safe for
\[
  r\le \left\lfloor \frac{B_Q-1}{d}\right\rfloor-1 .
\]
The first unsafe radius is \(r=\lfloor(B_Q-1)/d\rfloor\), if this radius is in
range.
\end{enumerate}
Negative right-hand sides mean that even radius zero is not certified for that
ledger and denominator.
\end{corollary}

\subsection{Prize-rate applicability}

For the Proximity Prize rates \(\rho\in\{1/2,1/4,1/8,1/16\}\), with
\(k\le2^{40}\), the exact line range reaches integer radii
\[
  r\le \frac{n-k}{3}=\frac{(1-\rho)n}{3}=\frac{1-\rho}{3\rho}k .
\]
Thus, for a field with \(|F|\approx2^\lambda\), the high-agreement tangent gate
can pin the line threshold only when roughly
\[
  2^{\lambda-128}\lesssim \frac{n-k}{3} .
\]
At the maximal allowed dimension \(k=2^{40}\), this gives the following
approximate largest field-bit sizes for which the line threshold can still be
pinned inside the exact tangent range:
\[
\begin{array}{c|c|c}
\rho & (n-k)/3\text{ at }k=2^{40} & \lambda\text{ up to about}\\
\hline
1/2  & 2^{40}/3      & 166.4\\
1/4  & 2^{40}        & 168.0\\
1/8  & 7\cdot2^{40}/3 & 169.2\\
1/16 & 5\cdot2^{40} & 170.3
\end{array}
\]
For substantially larger challenge fields, such as \(2^{192}\) or \(2^{256}\),
the tangent numerator is below the \(2^{-128}\) budget throughout the exact
high-agreement range for all prize-size rows.  The threshold, if any, must then
be controlled by the lower-agreement quotient-core, generated-field entropy,
curve/folding, or aperiodic local-limit ledgers; the tangent-star method alone
no longer pins it.
